\documentclass{controle}
\usepackage{main}

\title{Le paradoxe des bancs}
\author{Première Spécialité Mathématiques}
\date{Chapitre 6}

\begin{document}
\maketitle
\thispagestyle{empty}

\begin{tcolorbox}
Dans un square, il y a trois bancs à deux places. Mohammed et Natalie viennent s'asseoir successivement en choisissant une place au hasard. Mohammed s'assoit en premier, et Natalie ne peut pas choisir la même place que Mohammed. Quelles sont les chances que Mohammed et Natalie se retrouvent sur le même banc ? 
\end{tcolorbox}
\vspace*{0.2cm}

\begin{questions}
\titledquestion{Représenter l'expérience aléatoire}
\begin{parts}

\vspace*{0.5cm}
\begin{minipage}{0.4\textwidth}
\part Proposer un tableau à double entrée pour représenter l'ensemble des issues possibles.
\vspace*{0.2cm}

\makeemptybox{8cm}
\end{minipage}
\hfill\vline\hfill
\begin{minipage}{0.4\textwidth}
\part De la même manière, proposer un arbre de dénombrement pour représenter l'expérience.
\vspace*{0.5cm}

\makeemptybox{8cm}
\end{minipage}
\vspace*{0.2cm}
\part Comment utiliser ces représentations pour en déduire la probabilité recherchée ?
\part Comparer la réponse de la question précédente avec d'autres membres de la classe. Que remarquez-vous ?
\end{parts}
\end{questions}
\end{document}